\chapter{Products and coefficients}
\label{ch:products-coefficients}

The square of a primitive distinguished a differential graded algebra
from its cohomology algebra. Retaining only dimensions discards still
more information. We ask what the dimensions in each degree determine.

\section{Counting monomials}

Fix a field $k$ of characteristic zero for the algebraic counting
constructions. Write $\mathbb N=\{0,1,2,\ldots\}$.
The counting degree is a nonnegative integer weight, independent of any
cohomological degree.

The polynomial algebra $k[x]$ consists of finite linear combinations
of $x^a$, with $a\in\mathbb N$ and multiplication $x^ax^b=x^{a+b}$.
Give $x$ degree one. Its degree-$n$ subspace has basis $x^n$
and dimension one.
In $k[x,y]$, give both variables degree one.
The monomial $x^ay^b$ has degree $a+b$, so the degree-$n$ basis is
\[
 \{x^a y^{n-a}:a=0,1,\ldots,n\}.
\]
Its dimension is $n+1$.

These counts turn the choice of generators into a sequence of integers.
Recovering the generators from the sequence is a different question.
Cancellation already makes that question ambiguous when some generators
contribute with a minus sign.

The symbol $p$ can record these sequences as
\[
 1+p+p^2+p^3+\cdots,
 \qquad
 1+2p+3p^2+4p^3+\cdots.
\]
At this point each expression records coefficients, with no numerical value
assigned to $p$.

To count generators with signs, a second grading is needed.
A parity is an element of $\mathbb Z/2\mathbb Z=\{\bar0,\bar1\}$.
A vector space with degree and parity has a decomposition
\[
 V=\bigoplus_{n\geq0}(V_{n,\bar0}\oplus V_{n,\bar1}).
\]
Write $V_n=V_{n,\bar0}\oplus V_{n,\bar1}$.

The terms even and odd refer to parity. Counting degree and parity are
independent. Products add degrees in $\mathbb N$ and parities in
$\mathbb Z/2\mathbb Z$.
Scalars have degree zero and even parity, and algebra identities are even.

For a homogeneous element $a$, write $\epsilon(a)\in\{0,1\}$ for its
parity. Tensor products of algebras in this chapter use the multiplication
\[
 (a\otimes b)(a'\otimes b')
   =(-1)^{\epsilon(b)\epsilon(a')}aa'\otimes bb'.
\]
Thus parity determines the interchange sign. Counting degree contributes
no sign. On homogeneous tensors $a\otimes b$, $a'\otimes b'$, and
$a''\otimes b''$, both parenthesizations have exponent
\[
 \epsilon(b)\epsilon(a')+
 \epsilon(b)\epsilon(a'')+\epsilon(b')\epsilon(a'')
\pmod 2.
\]
Thus this multiplication is associative, with identity $1\otimes1$.

The exterior algebra on one generator $\theta$ is the vector space
$\Lambda(\theta)=k\,1\oplus k\theta$, with $1$ as identity and
$\theta^2=0$. Give $\theta$ counting degree three and odd parity, and give
$x,y$ counting degree one and even parity. Then
\[
 B=k[x,y]\otimes_k\Lambda(\theta),
 \qquad \deg x=\deg y=1,\quad\deg\theta=3,
\]
identifies as a vector space with $k[x,y]\oplus k[x,y]$ by
\[
 f\otimes(a+b\theta)\longmapsto(af,bf),\qquad a,b\in k.
\]
Its inverse sends $(f,g)$ to $f\otimes1+g\otimes\theta$.
Thus $B$ has basis $x^a y^b$ and $x^a y^b\theta$.
The tensor rule makes the polynomial generators commute with $\theta$.

When both spaces in degree $n$ are finite-dimensional, define
\[
 \dim V_n=\dim V_{n,\bar0}+\dim V_{n,\bar1},
 \qquad
 \operatorname{sdim}V_n=\dim V_{n,\bar0}-\dim V_{n,\bar1}.
\]
The second integer is the signed dimension, also called the superdimension.
Define the dimension series and the signed dimension series by
\[
 H_V(p)=\sum_{n\geq0}\dim V_n\,p^n,
 \qquad
 S_V(p)=\sum_{n\geq0}\operatorname{sdim}V_n\,p^n.
\]

In $B$, the monomials without $\theta$ are even, and those with $\theta$ are
odd. There are $n+1$ even monomials in degree $n$.
For $n\geq3$, there are $n-2$ odd monomials, because $a+b=n-3$.
For $n<3$, there are none. Thus
\begin{equation}
 S_B(p)=1+2p+3p^2+3p^3+3p^4+\cdots.
 \label{eq:pc-count-example}
\end{equation}

The ordinary dimension in degree $n\geq3$ is $2n-1$.
Consequently, even positive signed coefficients need not be ordinary
dimensions of the space from which they were computed.

\section{Formal multiplication and integer powers}

The choices of powers of different generators are independent.
Multiplying their counting series must therefore sum over all ways to split
a degree. This gives a multiplication defined entirely by finite sums.

Let $R$ be a commutative ring with identity. Its formal power-series ring
$R[[p]]$ consists of sequences $(a_0,a_1,\ldots)$, written
$\sum_{n\geq0}a_np^n$. Addition is coefficientwise, and multiplication is
\[
 \left(\sum_{n\geq0}a_np^n\right)
 \left(\sum_{n\geq0}b_np^n\right)
 =\sum_{n\geq0}\left(\sum_{j=0}^n a_jb_{n-j}\right)p^n.
\]
The notation $[p^n]F$ means the coefficient of $p^n$ in $F$.

Each inner sum is finite. Associativity follows because either bracketing
of three factors gives $\sum_{i+j+k=n}a_i b_j c_k$ in degree $n$.
Commutativity and distributivity follow from the corresponding laws in $R$.
The series with constant coefficient $1$ and all other coefficients zero
is the multiplicative identity.

For $N\geq0$, write $F\equiv G\pmod{p^{N+1}}$ when their coefficients
through degree $N$ agree. Equivalently,
$F-G\in p^{N+1}R[[p]]$. This congruence respects addition and multiplication.

A sequence of series converges formally when, in each bounded range of
degrees, its coefficients eventually stabilize. The limit is the sequence
of those stable coefficients. This is the $p$-adic topology used here.
It uses stabilization of coefficients, without a topology on $R$.

\begin{lemma}[Formal units]
\label{lem:pc-units}
A series $U=\sum_{n\geq0}u_np^n$ has a multiplicative inverse in $R[[p]]$
if and only if $u_0$ has a multiplicative inverse in $R$.
In particular, every series in $1+pR[[p]]$ is a unit.
\end{lemma}

\begin{proof}
If $UV=1$, the constant coefficients satisfy $u_0v_0=1$.
Conversely, suppose $u_0$ is a unit. Define
\[
 v_0=u_0^{-1},\qquad
 v_n=-u_0^{-1}\sum_{j=1}^n u_jv_{n-j}\quad(n\geq1).
\]
The coefficient formula for multiplication gives $[p^0]UV=1$ and
$[p^n]UV=0$ for every $n\geq1$. Thus $UV=1$.
The same recursion forces every coefficient of an inverse, so it is unique.
\end{proof}

For $j\geq1$, define the factorial $j!=1\cdot2\cdots j$, and set $0!=1$.
For an integer $m$ and $k\geq0$, define
\[
 \binom{m}{0}=1,\qquad
 \binom{m}{k}=\frac{m(m-1)\cdots(m-k+1)}{k!}\quad(k\geq1).
\]

For $0\leq k\leq m$, there are $m(m-1)\cdots(m-k+1)$ ordered choices
of $k$ distinct elements from an $m$-element set.
Each $k$-element subset has $k!$ orderings, so the quotient counts those
subsets. It is therefore an integer.
For $m\geq0$ and $k>m$, the numerator contains a zero factor.
Thus these coefficients vanish, as does the subset count.

For $m=-h<0$, the defining formula gives
\[
 (-1)^k\binom{-h}{k}=\binom{h+k-1}{k}.
\]
An integer coefficient in $R$ means its image under the map
$\mathbb Z\to R$ that sends $1$ to the identity of $R$.

\begin{lemma}[Binomial factors]
\label{lem:pc-binomial}
For $c\in R$, $\nu\geq1$, and $m\in\mathbb Z$, the factor
$(1-cp^\nu)^m$ is a unit of $R[[p]]$, and
\[
 (1-cp^\nu)^m
   =\sum_{k\geq0}(-1)^k\binom{m}{k}c^k p^{k\nu}.
\]
Here a negative power means the corresponding positive power of the
inverse from Lemma~\ref{lem:pc-units}.
\end{lemma}

\begin{proof}
First use a formal variable $t$. For $m\geq0$, expand the product of $m$
copies of $1-t$. A term $(-t)^k$ occurs once for each choice of $k$ copies,
which gives the formula.
The coefficient rule also gives
\[
 (1-t)\sum_{k\geq0}t^k=1.
\]

For $h\geq1$, the coefficient of $t^k$ in
$(\sum_{j\geq0}t^j)^h$ counts tuples
$(a_1,\ldots,a_h)\in\mathbb N^h$ with sum $k$.
Place $k$ identical marks in a row and insert $h-1$ separators.
Empty groups are allowed. This identifies the tuples with choices of
the separator positions among $k+h-1$ positions, giving
$\binom{k+h-1}{h-1}=\binom{k+h-1}{k}$ tuples.
This proves the formula for $m=-h$ in $\mathbb Z[[t]]$.

Send integer coefficients to $R$ and substitute $t=cp^\nu$.
The substitution is defined coefficientwise because $\nu\geq1$.
All factors have constant coefficient one, so they are units.
\end{proof}

The one-generator counting sequence is therefore $(1-p)^{-1}$.
Multiplying two copies counts the monomials in $k[x,y]$.
The two signed contributions of $1$ and $\theta$ give $1-p^3$.
Thus the count in \eqref{eq:pc-count-example} also gives
\begin{equation}
 S_B(p)=\frac{1-p^3}{(1-p)^2}.
 \label{eq:pc-rational-example}
\end{equation}
The fraction means multiplication by a formal inverse, whose constant
coefficient is one.

\section{Locally finite products and weighted counts}

Infinitely many generators are compatible with finite coefficients when
only finitely many generators can contribute in each bounded degree.
The same condition defines an infinite product without choosing analytic
values of $p$.

Let $I$ be an index set. For each $r\in I$, specify a positive integer
$\nu_r$, an integer $m_r$, and an element $c_r\in R$.
Assume that
\begin{equation}
 I_{\leq N}=\{r\in I:\nu_r\leq N\}\quad\text{is finite for every }N.
 \label{eq:pc-local-finiteness}
\end{equation}
This is the local finiteness condition on the degrees.

\begin{proposition}[Coefficientwise product]
\label{prop:pc-product}
Under \eqref{eq:pc-local-finiteness}, there is a unique series
\begin{equation}
 F(p)=\prod_{r\in I}(1-c_rp^{\nu_r})^{m_r}\in1+pR[[p]]
 \label{eq:pc-product}
\end{equation}
whose coefficients through degree $N$ equal those of the finite product
over $I_{\leq N}$, for every $N$.
The result is independent of the order of the factors.
\end{proposition}

\begin{proof}
Lemma~\ref{lem:pc-binomial} shows that the $r$th factor is congruent to
$1$ modulo $p^{\nu_r}$. A factor with $\nu_r>N$ therefore leaves every
coefficient through degree $N$ unchanged.
The finite product over $I_{\leq N+1}$ consequently agrees with the product
over $I_{\leq N}$ through degree $N$.
These compatible coefficients define one series. They also force uniqueness.
Any reordering gives the same finite products modulo each $p^{N+1}$,
because multiplication in $R[[p]]$ is commutative.
\end{proof}

Positive degrees alone do not imply local finiteness. For example,
the partial products $(1-p)^k$ have coefficient $-k$ in degree one
over $\mathbb Z$. This coefficient never stabilizes, so the repeated
infinite product of the factor $1-p$ has no formal limit.
By contrast, $\prod_{n\geq1}(1-p^n)$ exists by
Proposition~\ref{prop:pc-product}.
Degree zero also requires separate treatment: $1-cp^0=1-c$ need not
be a unit.

The integer exponents in \eqref{eq:pc-product} have a direct counting
interpretation. For each $r$ with $m_r<0$, introduce $-m_r$ even polynomial
generators of degree $\nu_r$. For each $r$ with $m_r>0$, introduce $m_r$
odd exterior generators of that degree. Indices with $m_r=0$ introduce none.
Let $V_{\bar0}$ and $V_{\bar1}$ be the $k$-vector spaces on the respective
sets of generators. Define
\[
 \mathcal A=\operatorname{Sym}(V_{\bar0})\otimes_k
                 \Lambda(V_{\bar1}).
\]

Here $\operatorname{Sym}(V_{\bar0})$ is the polynomial algebra with finite
monomials in the even generators. The exterior algebra
$\Lambda(V_{\bar1})$ has basis the finite products of distinct odd
generators in a fixed order.
Interchanging two odd generators changes the sign, and the square of each
odd generator is zero. Even generators commute with every generator.
These rules define the multiplication on the displayed basis.

Associativity follows by counting interchanges among three ordered factors.
If an odd generator repeats, both parenthesizations are zero. Otherwise,
both signs count the same pairs of odd generators whose order is reversed.
For homogeneous elements $a,b$ with parities $\epsilon,\delta\in\{0,1\}$,
these rules give $ab=(-1)^{\epsilon\delta}ba$.

Weighted counting assigns a scalar to each monomial.
For a linear operator $T:V\to V$, a nonzero vector $v$ satisfying
$T(v)=\lambda v$ is an eigenvector, with eigenvalue $\lambda\in k$.
In a basis of eigenvectors, the matrix has these eigenvalues on its
diagonal and zero in every other entry.
The trace calculation also requires comparison with an arbitrary basis.

For a finite-dimensional vector space $V$, let
$\mathcal B=(v_1,\ldots,v_d)$ and $\mathcal C=(w_1,\ldots,w_d)$ be ordered
bases. Write $[v]_{\mathcal B}$ and $[v]_{\mathcal C}$ for coordinate
columns. Let $S$ have $j$th column $[w_j]_{\mathcal B}$.
Then $[v]_{\mathcal B}=S[v]_{\mathcal C}$.
The reverse coordinate change is an inverse for $S$.

If $T$ has matrices $M_{\mathcal B}$ and $M_{\mathcal C}$ in these bases,
applying $T$ to the coordinate identity gives
$M_{\mathcal B}S=SM_{\mathcal C}$.
Hence $M_{\mathcal C}=S^{-1}M_{\mathcal B}S$.
For a finite-dimensional space, the trace of an operator is the sum of
the diagonal entries of any matrix representing it.
This does not depend on the basis: finite sums give
$\operatorname{tr}(AB)=\sum_{i,j}A_{ij}B_{ji}=\operatorname{tr}(BA)$,
and hence $\operatorname{tr}(S^{-1}AS)=\operatorname{tr}(A)$.
For an operator that preserves parity, its supertrace is the trace on the
even subspace minus the trace on the odd subspace.
When comparing traces with dimension counts, regard integers as their
images under the injective scalar map $\mathbb Z\to k$.

\begin{proposition}[Signed monomial realization]
\label{prop:pc-realization}
Assume \eqref{eq:pc-local-finiteness} and take $R=k$.
Each degree of $\mathcal A$ is finite-dimensional.
Define the linear operator $T$ by multiplying each generator of index $r$
by $c_r$, and extend $T$ multiplicatively with $T(1)=1$.
Then $T$ preserves degree and parity, and
\[
 \sum_{n\geq0}\operatorname{str}(T|_{\mathcal A_n})p^n
   =\prod_{r\in I}(1-c_rp^{\nu_r})^{m_r}.
\]
In particular, if all $c_r=1$, the product is $S_{\mathcal A}(p)$.
\end{proposition}

\begin{proof}
A degree-$n$ monomial uses only generators of degree at most $n$.
There are finitely many such generators, and every polynomial exponent is
at most $n$. This proves finite-dimensionality, including
$\mathcal A_0=k$.
Scaling generators respects every multiplication relation, so $T$ is
well-defined. Each monomial is an eigenvector, with eigenvalue the product
of its generator weights.

For one even generator of degree $\nu$ and weight $c$, its contribution is
$\sum_{a\geq0}c^a p^{a\nu}=(1-cp^\nu)^{-1}$.
For one odd generator, the basis $1,\theta$ contributes $1-cp^\nu$.

In a product monomial, eigenvalues multiply, degrees add, and parity signs
multiply. The coefficient rule therefore multiplies these single-generator
series. In each degree this argument uses finitely many generators.
Proposition~\ref{prop:pc-product} then gives the stated identity.
\end{proof}

\section{Logarithmic coefficients and the recurrence}

Differentiating a product and dividing by that product gives a sum of
contributions from its individual factors. This identity yields a
coefficient recurrence.
For $F=\sum_{n\geq0}f_np^n\in R[[p]]$, define its formal derivative by
\[
 F'=\sum_{n\geq1}n f_np^{n-1}.
\]
The coefficient rule proves $(FG)'=F'G+FG'$:
the coefficient with $i+j=n$ has multiplier $n=i+j$ on both sides.
If $F$ is a unit, differentiation of $FF^{-1}=1$ gives
$(F^{-1})'=-F^{-2}F'$.
Thus the usual integer-power rule holds for positive, zero, and negative
powers of units.

For $F\in1+pR[[p]]$, define its logarithmic coefficients, or ghost
coefficients, by
\begin{equation}
 \frac{pF'(p)}{F(p)}=\sum_{s\geq1}A_s p^s.
 \label{eq:pc-ghost-definition}
\end{equation}
The quotient uses the formal inverse of $F$.
This definition requires no logarithm and no division by integers in $R$.

\begin{theorem}[Product recurrence]
\label{thm:pc-recurrence}
For the product in \eqref{eq:pc-product}, its logarithmic coefficients are
\begin{equation}
 A_s=-\sum_{\substack{r\in I\\\nu_r\mid s}}
                  m_r\nu_r c_r^{s/\nu_r}\qquad(s\geq1).
 \label{eq:pc-ghost-product}
\end{equation}
Writing $F=\sum_{j\geq0}f_jp^j$, one has $f_0=1$ and
\begin{equation}
 j f_j=\sum_{s=1}^j A_s f_{j-s}\qquad(j\geq1).
 \label{eq:pc-recurrence}
\end{equation}
Here $\nu_r\mid s$ means that $s$ is an integer multiple of $\nu_r$.
Each sum in \eqref{eq:pc-ghost-product} is finite.
\end{theorem}

\begin{proof}
For a single factor, the derivative and geometric-series formulas give
\[
 \frac{p\bigl((1-cp^\nu)^m\bigr)'}{(1-cp^\nu)^m}
   =-\frac{m\nu cp^\nu}{1-cp^\nu}
   =-m\nu\sum_{k\geq1}c^k p^{k\nu}.
\]

For a finite product, the product rule adds these expressions.
For degree at most $N$, the infinite product can be replaced by the
finite product over $I_{\leq N}$.
Indeed, $p$ times differentiation preserves congruence modulo $p^{N+1}$.
Inversion also preserves that congruence for units, since
$U^{-1}-V^{-1}=U^{-1}(V-U)V^{-1}$.
Therefore the same additive formula holds in every degree for $F$.

The coefficient of $p^s$ in the single-factor expression is nonzero only
when $\nu\mid s$, which proves \eqref{eq:pc-ghost-product}.
Only indices in the finite set $I_{\leq s}$ can occur.
Multiplying \eqref{eq:pc-ghost-definition} by $F$ and comparing the
coefficient of $p^j$ proves \eqref{eq:pc-recurrence}.
\end{proof}

For a commutative $\mathbb Q$-algebra $R$, every positive integer is
invertible in $R$. Thus \eqref{eq:pc-recurrence} recovers $f_j$ uniquely
from $A_1,\ldots,A_j$ and the preceding coefficients.
The converse recursion is
\[
 A_j=jf_j-\sum_{s=1}^{j-1}A_s f_{j-s}.
\]
Hence the first $N$ logarithmic coefficients contain exactly the same
information as $f_1,\ldots,f_N$ over such a ring.

The exponential form of this correspondence can also be constructed
coefficientwise. For $H\in pR[[p]]$, define
\[
 \exp(H)=\sum_{k\geq0}\frac{H^k}{k!}.
\]
Since $H^k\in p^kR[[p]]$, each coefficient is a finite sum.
The power rule gives $(\exp H)'=H'\exp H$ by differentiating those finite
sums degree by degree. Consequently, for any sequence $A_s\in R$,
\begin{equation}
 F(p)=\exp\left(\sum_{s\geq1}\frac{A_s}{s}p^s\right)
 \label{eq:pc-exponential}
\end{equation}
has constant coefficient one and satisfies
\eqref{eq:pc-ghost-definition}. The recurrence proves uniqueness.

For example, integers $A_1=1$ and $A_s=0$ for $s>1$ give
$F=\exp(p)=1+p+\tfrac12p^2+\cdots$. Integer logarithmic coefficients
therefore do not by themselves give integer coefficients of $F$.

If a family of local factors is specified by
\[
 E_r(X)=\exp\left(\sum_{k\geq1}\frac{a_{r,k}}{k}X^k\right)
 \in1+X\mathbb Q[[X]],
\]
and positive integers $d_r$ satisfy the local finiteness condition, then
$\prod_r E_r(p^{d_r})$ exists formally by the same truncation argument.
Here local finiteness means that only finitely many $r$ have $d_r\leq N$,
for each $N$. Applying the product rule gives its logarithmic coefficients
\[
 A_s=\sum_{d_r\mid s}d_r a_{r,s/d_r}.
\]
Thus this construction also has recurrence \eqref{eq:pc-recurrence}.

\section{Integer coefficients and factors of weight one}

An exact criterion distinguishes the rational sequences $(A_s)$ for
which \eqref{eq:pc-exponential} has integer coefficients.
Factors with weight one give this criterion and also a unique
factorization of every series in $1+p\mathbb Z[[p]]$.
The sign convention in this section is
\[
 F(p)=\prod_{n\geq1}(1-p^n)^{-b_n},
\]
so $b_n$ is the negative of the exponent $m_n$ used earlier.

Divisors of an integer will relate its logarithmic coefficient to the
factor exponents. The required divisor identity uses unique factorization
into primes, which follows from integer division as follows.
For positive integers $a,d$, the set of nonnegative integers $q$ with
$qd\leq a$ is nonempty and finite. Its largest element satisfies
\[
 a=qd+r,\qquad 0\leq r<d.
\]
The upper bound on $r$ follows because $(q+1)d>a$.

A prime is an integer $\ell>1$ whose positive divisors are $1$ and $\ell$.
If a prime $\ell$ divides $ab$ for positive integers $a,b$, then it
divides $a$ or $b$. To prove this, suppose that $\ell$ does not divide $a$.
Let $d$ be the least positive integer of the form $u\ell+va$ with
$u,v\in\mathbb Z$. Such a $d$ exists because $\ell$ is one such integer.

Divide both $\ell$ and $a$ by $d$ using the preceding construction.
Each remainder is an integer combination of $\ell$ and $a$.
If positive, it contradicts the minimality of $d$.
Thus $d$ divides both $\ell$ and $a$.

Primality and $\ell\nmid a$ force $d=1$.
For integers $u,v$ with $u\ell+va=1$, multiplication by $b$ gives
\[
 b=u\ell b+vab.
\]
Both terms on the right are divisible by $\ell$, so $\ell$ divides $b$.

Every integer $n>1$ is a finite product of primes. Indeed, induction on
$n$ proves this if $n$ is prime, and otherwise writes $n=ab$ with
$1<a,b<n$ and applies the induction hypothesis to both factors.

If two prime products are equal, the first prime in one divides a prime
in the other, by repeated application of the preceding divisibility result.
Those two primes are equal. Cancel them and repeat.
This proves uniqueness up to reordering. The number $1$ has the empty
prime factorization.

For a positive integer $d$, define its M\"obius function $\mu(d)$ as follows.
Set $\mu(1)=1$. Set $\mu(d)=0$ if a prime square divides $d$.
Otherwise $d$ is a product of $t$ distinct primes, and set $\mu(d)=(-1)^t$.

For any positive integer $q$,
\begin{equation}
 \sum_{d\mid q}\mu(d)=
 \begin{cases}1,&q=1,\\0,&q>1.\end{cases}
 \label{eq:pc-mobius-sum}
\end{equation}
To prove this, let $q>1$ have distinct prime divisors $\ell_1,\ldots,\ell_t$.
Only products of subsets of these primes contribute.
The sum is therefore $\sum_{J\subseteq\{1,\ldots,t\}}(-1)^{|J|}
=(1-1)^t=0$. For $q=1$, its only divisor is $1$.

\begin{lemma}[Divisor inversion]
\label{lem:pc-inversion}
For sequences $u_n,v_n$ in an abelian group, the identities
$v_n=\sum_{d\mid n}u_d$ for all $n\geq1$ hold if and only if
\[
 u_n=\sum_{d\mid n}\mu(d)v_{n/d}\quad(n\geq1).
\]
Integer multiplication here means repeated addition, with negatives taken
in the abelian group.
\end{lemma}

\begin{proof}
If $v_n=\sum_{d\mid n}u_d$, substitute this expression and rearrange the
finite divisor sums:
\[
 \sum_{d\mid n}\mu(d)v_{n/d}
   =\sum_{e\mid n}u_e\sum_{d\mid n/e}\mu(d)=u_n.
\]
The last equality is \eqref{eq:pc-mobius-sum}.
Conversely, substitute the claimed expression for $u_e$ to obtain
\[
 \sum_{e\mid n}u_e
   =\sum_{e\mid n}\sum_{d\mid e}\mu(d)v_{e/d}
   =\sum_{h\mid n}v_h\sum_{d\mid n/h}\mu(d)=v_n.
\]
All sums are finite, so both rearrangements are valid in an abelian group.
\end{proof}

\begin{theorem}[Integral factorization]
\label{thm:pc-integrality}
Every $F\in1+p\mathbb Z[[p]]$ has a unique expression
\[
 F=\prod_{n\geq1}(1-p^n)^{-b_n},\qquad b_n\in\mathbb Z.
\]
More generally, for a sequence $A_s\in\mathbb Q$, let $F$ be
\eqref{eq:pc-exponential} in $\mathbb Q[[p]]$, and define
\begin{equation}
 b_n=\frac1n\sum_{d\mid n}\mu(d)A_{n/d}.
 \label{eq:pc-integrality-criterion}
\end{equation}
Then $F\in1+p\mathbb Z[[p]]$ if and only if $b_n\in\mathbb Z$
for every $n\geq1$.
\end{theorem}

\begin{proof}
First suppose $F$ has integer coefficients. Construct finite products
$P_n=\prod_{k=1}^n(1-p^k)^{-b_k}$ inductively, starting with $P_0=1$.
Suppose $P_{n-1}\equiv F\pmod{p^n}$.
Its inverse has integer coefficients by Lemma~\ref{lem:pc-units}, so
\[
 F P_{n-1}^{-1}\equiv1+b_n p^n\pmod{p^{n+1}}
\]
for a unique integer $b_n$.
Lemma~\ref{lem:pc-binomial} gives
$(1-p^n)^{-b_n}\equiv1+b_n p^n\pmod{p^{n+1}}$.
Thus $P_n\equiv F\pmod{p^{n+1}}$, completing the induction.

Proposition~\ref{prop:pc-product} defines the infinite product, and the
congruences show that it equals $F$.
Any factorization with integer exponents must choose this same $b_n$
after $b_1,\ldots,b_{n-1}$ have been fixed. This proves uniqueness.

The product formula \eqref{eq:pc-ghost-product}, with $m_n=-b_n$ and
$c_n=1$, gives
\[
 A_s=\sum_{n\mid s}n b_n.
\]
Apply Lemma~\ref{lem:pc-inversion} to $u_n=n b_n$.
It yields \eqref{eq:pc-integrality-criterion}. Thus integral $F$ forces
every number in that formula to be integral.

Conversely, suppose every number in \eqref{eq:pc-integrality-criterion}
is integral. Construct $G=\prod_{n\geq1}(1-p^n)^{-b_n}$ over $\mathbb Z$.
Divisor inversion gives $\sum_{n\mid s}n b_n=A_s$.
Hence $G$ has the prescribed logarithmic coefficients.
Viewed in $\mathbb Q[[p]]$, both $G$ and $F$ satisfy the same recurrence
with constant coefficient one. Its uniqueness gives $F=G$.
\end{proof}

For the sequence with $A_1=1$ and all later terms zero,
formula \eqref{eq:pc-integrality-criterion} gives $b_2=-\tfrac12$.
This recovers the obstruction already visible in $[p^2]\exp(p)=\tfrac12$.
The factorization uses coefficients in every degree.
A finite list requires an additional condition to determine the series.

\section{How many coefficients determine a series?}

A finite list of coefficients is a finite jet.
More precisely, the jet through degree $N$ of $F\in R[[p]]$ is its class
in $R[[p]]/p^{N+1}R[[p]]$, or equivalently the list
$([p^0]F,\ldots,[p^N]F)$.
For a field $K$, this list never determines an unrestricted formal series.
Indeed, $F$ and $F+p^{N+1}$ are distinct and have the same jet through
degree $N$. If $F(0)=1$, both remain normalized units.

The degree of a nonzero polynomial is its highest exponent with nonzero
coefficient. The product of two nonzero polynomials over $K$ is nonzero:
its highest coefficient is the product of their highest coefficients,
which is nonzero in a field.

A rational fraction regular at zero is a pair of polynomials $(N,D)$
with $D(0)\ne0$, written $N/D$.
Define equality of such fractions by
\[
 \frac ND=\frac{N'}{D'}
 \quad\Longleftrightarrow\quad ND'=N'D.
\]

This relation is reflexive and symmetric.
If $ND'=N'D$ and $N'D''=N''D'$, multiplication gives
$D'(ND''-N''D)=0$.
The product property just proved and $D'\ne0$ imply $ND''=N''D$.
Thus the relation is transitive.
A rational function regular at zero is an equivalence class for this
relation.

The formal expansion of $N/D$ is $ND^{-1}\in K[[p]]$.
Equality of two fractions is equivalent to equality of their formal
expansions: multiply by the formal unit $DD'$ in either direction.
Thus the expansion is well-defined and distinguishes rational functions.
Rescaling both polynomials by $D(0)^{-1}$ gives a representative with
$D(0)=1$.
Even knowledge that two series are rational is insufficient without a
degree bound: $1$ and $1+p^{N+1}$ give a counterexample.
In the next theorem, all four polynomials have constant coefficient one,
so their degrees are defined.

\begin{theorem}[Finite determination of rational functions]
\label{thm:pc-rational-determination}
Let
\[
 F=N_F D_F^{-1},\qquad G=N_G D_G^{-1},
\]
where $N_F,D_F,N_G,D_G\in K[p]$ all have constant coefficient one.
Set
\[
 M=\max\{\deg N_F+\deg D_G,\ \deg N_G+\deg D_F\}.
\]
If $F\equiv G\pmod{p^{M+1}}$, then $F=G$ as formal series and as
rational functions. In particular, normalized rational functions with
numerator degree at most $a$ and denominator degree at most $b$ are
determined by their coefficients through degree $a+b$.
\end{theorem}

\begin{proof}
Multiplication by $D_FD_G$ preserves the congruence, so the polynomial
\[
 H=N_FD_G-N_GD_F
\]
has all coefficients through degree $M$ equal to zero.
Both summands have degree at most $M$.
Therefore $H=0$.
Cross multiplication gives equality of rational functions.
Multiplication by the formal inverse of $D_FD_G$ gives equality of series.
If both numerator degrees are at most $a$ and both denominator degrees
are at most $b$, then $M\leq a+b$, giving the stated uniform bound.
\end{proof}

Every finite product \eqref{eq:pc-product} over a field is rational.
Explicitly,
\[
 N=\prod_{m_r>0}(1-c_rp^{\nu_r})^{m_r},\qquad
 D=\prod_{m_r<0}(1-c_rp^{\nu_r})^{-m_r}.
\]
Both have constant coefficient one, and
\[
 \deg N\leq\sum_{m_r>0}m_r\nu_r,\qquad
 \deg D\leq\sum_{m_r<0}(-m_r)\nu_r.
\]

Each polynomial $1-c_rp^{\nu_r}$ has degree at most $\nu_r$.
Multiplying the required copies proves these bounds.
They become equalities when all relevant $c_r$ are
nonzero. Cancellation between $N$ and $D$ can give smaller bounds for the
same rational function.

Without a bound on the degrees of the factors, no finite jet determines
a weight-one product. Given a normalized product $F$ and $N\geq0$,
the product $F(1-p^{N+1})$ has the same jet through degree $N$.
It is distinct from $F$, since $F$ is a unit and $1-p^{N+1}\ne1$.
Thus local finiteness and integral exponents alone do not supply a
finite bound for determination.

\section{The three-generator example revisited}

For $B=k[x,y]\otimes_k\Lambda(\theta)$, the two even degree-one generators
contribute $(1-p)^{-2}$ to \eqref{eq:pc-rational-example}.
The odd degree-three generator contributes $1-p^3$.
Formula \eqref{eq:pc-ghost-product} gives
\[
 A_s=2-3\,\mathbf1_{3\mid s},
\]
where $\mathbf1_{3\mid s}$ equals $1$ when $3$ divides $s$, and equals $0$
otherwise. Thus
\[
 A_1=2,\quad A_2=2,\quad A_3=-1,\quad A_4=2,\quad A_5=2,\quad A_6=-1.
\]

The recurrence starts with
\begin{align*}
 f_0&=1,& f_1&=2,\\
 2f_2&=2\cdot2+2\cdot1=6,&
 3f_3&=2\cdot3+2\cdot2-1\cdot1=9,\\
 4f_4&=2\cdot3+2\cdot3-1\cdot2+2\cdot1=12.
\end{align*}

For all degrees at once, use the polynomial identity
$1-p^3=(1-p)(1+p+p^2)$ to obtain
\[
 F=\frac{1+p+p^2}{1-p}.
\]
Multiplication by the geometric series gives $f_0=1$, $f_1=2$, and
$f_n=3$ for all $n\geq2$. The product recurrence applies to this same
formal series, and its uniqueness proves that the recurrence gives these
values in every degree.

The reduced expression has numerator degree two and denominator degree
one. By Theorem~\ref{thm:pc-rational-determination}, its coefficients
through degree three determine it among normalized rational functions
with those bounds. In the larger class with numerator degree at most
three and denominator degree at most two, coefficients through degree
five suffice. Neither statement identifies $F$ among arbitrary formal
series from the same finite data.

The same coefficients also count an ordinary graded algebra:
\[
 C=k[x,y]/(x^3),\qquad\deg x=\deg y=1.
\]
Here $(x^3)$ is the ideal of multiples of $x^3$.
It is spanned by the monomials $x^a y^b$ with $a\geq3$.
Consequently the residue classes with $0\leq a\leq2$ form a basis of $C$.
There are $\min(n+1,3)$ such monomials in degree $n$.

With every vector declared even, this gives $H_C=S_C=S_B$.
Yet $B$ has odd vectors and has ordinary dimension $2n-1$ in degree
$n\geq3$. Equality of the scalar series has forgotten that distinction.

\section{Lattice labels at a fixed frame}

Lattice data can index a product whose weights depend on a chosen frame.
The following example gives the three-generator product at one frame and
changes its degree-one coefficient under a quarter-turn.

A free abelian group $L$ of rank $d$ has elements $v_1,\ldots,v_d$ such that
each element has a unique expression $\sum_{j=1}^d a_jv_j$ with
$a_j\in\mathbb Z$.
An integer-valued symmetric bilinear form is a map
$L\times L\to\mathbb Z$ additive in each argument and unchanged by
interchanging its arguments.
A lattice here means such a group together with such a form.
Take
\[
 L=\mathbb Z^3,\qquad
 (v,w)=v_1w_1+v_2w_2-2v_3w_3.
\]

The norm of $v$ means $(v,v)$, which can be negative.
Use the field $\mathbb R$ of real numbers with its total order.
The order satisfies $x\leq y\Rightarrow x+z\leq y+z$ and
$0\leq x,\ 0\leq y\Rightarrow0\leq xy$.
Extend the displayed formula to real coordinates in
$L_{\mathbb R}=\mathbb R^3$.
The standard basis vectors are $e_1=(1,0,0)$, $e_2=(0,1,0)$, and
$e_3=(0,0,1)$. The form has matrix $\operatorname{diag}(1,1,-2)$ in
this basis.

The plane
$P=\operatorname{span}_{\mathbb R}(e_1,e_2)=\{(x,y,0):x,y\in\mathbb R\}$
is positive: the norm of a nonzero vector in $P$ is $x^2+y^2>0$.
A frame of $P$ is an ordered basis $(f_1,f_2)$.
It is orthonormal if $(f_j,f_j)=1$ for $j=1,2$ and $(f_1,f_2)=0$.

For a real matrix $M=\left(\begin{smallmatrix}a&b\\c&d\end{smallmatrix}\right)$,
define $\det M=ad-bc$.
For $N=\left(\begin{smallmatrix}e&f\\g&h\end{smallmatrix}\right)$, expansion gives
\[
 \begin{aligned}
 \det(MN)
   &=(ae+bg)(cf+dh)-(af+bh)(ce+dg)\\
   &=(ad-bc)(eh-fg)=\det(M)\det(N).
 \end{aligned}
\]
For an invertible $M$, this implies $\det M\ne0$ and
$\det(M^{-1})=1/\det M$.
Two frames have the same orientation when the matrix changing one to the
other has positive determinant.
The product and inverse identities make this an equivalence relation.

Fix the orientation containing $(e_1,e_2)$.
The displayed bilinear form gives $(e_1,e_1)=(e_2,e_2)=1$ and
$(e_1,e_2)=0$, so this frame is orthonormal.
Write $Z=(P,e_1,e_2)$ for this fixed datum.

\subsection{Numerical exponential weights}

The period coordinate will take complex values, and its weight will be a
numerical exponential. These values require convergence in the scalar
field. The variable $p$ remains formal.

A set $A\subseteq\mathbb R$ is bounded above if some $M\in\mathbb R$
satisfies $a\leq M$ for every $a\in A$.
Such an $M$ is an upper bound.
Completeness of $\mathbb R$ means that every nonempty set bounded above
has a least upper bound, called its supremum.
For a nonempty set bounded below, define its infimum by
$\inf A=-\sup(-A)$.

The positive integers are unbounded.
Indeed, if their supremum were $s$, an integer $n>s-1$ would give
$n+1>s$. Since $2^n\geq n+1$, the number $2^{-n}$ can be made smaller
than any positive real number.
For a real number $x$, put $|x|=\max(x,-x)$.

Every $r\geq0$ has a unique nonnegative square root.
For $r>0$, take
\[
 s=\sup\{u\geq0:u^2\leq r\}.
\]
The set contains $\min(1,r)/2>0$ and has upper bound $r+1$.
If $s^2<r$, choose
$0<\delta<\min(1,(r-s^2)/(2s+1))$.
Then $(s+\delta)^2<r$, contradicting the upper bound $s$.

If $s^2>r$, choose
$0<\delta<\min(s,(s^2-r)/(2s))$.
Then $(s-\delta)^2>r$, so $s-\delta$ is already an upper bound for the
same set. This contradicts the definition of $s$.
Thus $s^2=r$.

Uniqueness follows because squaring is strictly increasing on
nonnegative numbers. The square root of zero is zero.

Define $\mathbb C$ to consist of pairs of real numbers, written $x+iy$,
with operations
\[
 \begin{aligned}
 (x+iy)+(u+iv)&=(x+u)+i(y+v),\\
 (x+iy)(u+iv)&=(xu-yv)+i(xv+yu).
 \end{aligned}
\]
Real numbers identify with pairs $x+i0$, and $i=0+i1$ satisfies $i^2=-1$.
Addition has identity $0+i0$ and inverse $-(x+iy)=-x-i y$.
Multiplication is commutative, has identity $1+i0$, and distributes over
addition by expansion.
For a third factor $a+ib$, both parenthesizations of
$(x+iy)(u+iv)(a+ib)$ give
\[
 (xua-yva-xvb-yub)+i(xub-yvb+xva+yua).
\]
This proves associativity.

A nonzero $z=x+iy$ has inverse $(x-iy)/(x^2+y^2)$.
Thus $\mathbb C$ is a field, and its real subfield makes its
characteristic zero.
Write
\[
 \operatorname{Re}z=x,\qquad \operatorname{Im}z=y,\qquad
 \overline z=x-iy,\qquad |z|=\sqrt{x^2+y^2}.
\]
These are its real part, imaginary part, conjugate, and modulus,
respectively.

Expansion gives $|zw|=|z||w|$ and
\[
 (x^2+y^2)(u^2+v^2)-(xu+yv)^2=(xv-yu)^2\geq0.
\]
The latter identity gives $xu+yv\leq|z||w|$ and hence
$|z+w|\leq|z|+|w|$ by squaring both nonnegative sides.
Applying this triangle inequality to $z=(z-w)+w$ and then interchanging
$z,w$ gives $\bigl||z|-|w|\bigr|\leq|z-w|$.

A sequence $z_n$ converges to $z$ if, for every $\varepsilon>0$,
there is $N$ such that $|z_n-z|<\varepsilon$ for all $n\geq N$.
It is Cauchy if, for every $\varepsilon>0$, there is $N$ such that
$|z_n-z_m|<\varepsilon$ for all $m,n\geq N$.
Every real Cauchy sequence $a_n$ is bounded: its tail lies within one
of a fixed term, and its initial segment is finite.
Put
\[
 l_N=\inf_{n\geq N}a_n,\qquad u_N=\sup_{n\geq N}a_n,\qquad
 a=\sup_N l_N.
\]
For every $M,N$, any $a_n$ with $n\geq\max(M,N)$ gives
$l_M\leq a_n\leq u_N$. Thus $l_N\leq a\leq u_N$.

Given $\varepsilon>0$, the Cauchy condition makes
$a_m\leq a_n+\varepsilon$ for all $m,n\geq N$, once $N$ is large enough.
Taking the supremum in $m$ and the infimum in $n$ gives
$u_N-l_N\leq\varepsilon$.
Since both $a_n$ and $a$ lie in $[l_N,u_N]$ for $n\geq N$, this proves
$a_n\to a$.
The inequalities
\[
 \max(|x|,|y|)\leq|x+iy|\leq|x|+|y|
\]
show that a complex sequence is Cauchy, or convergent, exactly when both
coordinates have that property. Hence every complex Cauchy sequence
converges.

The triangle inequality makes limits unique and proves the limit rule
for sums. Real limits preserve weak inequalities by the defining
$\varepsilon$ bound.
For products, boundedness of a convergent sequence and the estimate
\[
 |z_nw_n-zw|
 \leq |z_n|\,|w_n-w|+|w|\,|z_n-z|
\]
give the product rule for limits.
Conjugation and the two coordinate maps also preserve limits.
The finite identity
$(1-q)\sum_{j=0}^m q^j=1-q^{m+1}$ supplies the geometric bounds below.

For $z\in\mathbb C$, define the finite sums
\[
 E_N(z)=\sum_{n=0}^N\frac{z^n}{n!}.
\]
For fixed $r\geq0$, the terms $a_n=r^n/n!$ satisfy
$a_{n+1}=ra_n/(n+1)\leq a_n/2$ once $n+1\geq2r$.
Thus $a_n\to0$, and for all sufficiently large $N$ and every $M>N$,
\[
 \sum_{n=N+1}^M\frac{r^n}{n!}
 \leq\frac{2r^{N+1}}{(N+1)!}\longrightarrow0.
\]
With $r=|z|$, this estimate proves that $E_N(z)$ is Cauchy.
It defines the numerical exponential on all of $\mathbb C$:
\[
 \exp z=\lim_{N\to\infty}E_N(z)
       =\sum_{n\geq0}\frac{z^n}{n!}.
\]
The same estimate proves absolute convergence, meaning convergence of
$\sum_{n\geq0}|z^n/n!|$.

Expanding $n$ copies of $z+w$ and grouping the choices of $z$ gives
$(z+w)^n=\sum_{j=0}^n\binom nj z^j w^{n-j}$.
This identifies the terms of total degree at most $N$ in
$E_N(z)E_N(w)$ with $E_N(z+w)$.
Every remaining term has total degree between $N+1$ and $2N$, so
\[
 |E_N(z)E_N(w)-E_N(z+w)|
 \leq\sum_{n=N+1}^{2N}\frac{(|z|+|w|)^n}{n!}\longrightarrow0.
\]
Taking limits gives
\[
 \exp(z+w)=\exp z\,\exp w,\qquad
 \exp0=1,\qquad
 \exp\overline z=\overline{\exp z}.
\]
In particular, $\exp z\,\exp(-z)=1$.

For real $x$, $\exp x$ is real and
$\exp x=(\exp(x/2))^2>0$.
For $x>0$, the nonnegative terms of its series give
$\exp x\geq1+x>1$.
For $x<0$, the identity $\exp x=1/\exp(-x)$ gives $0<\exp x<1$.
Consequently
\[
 \begin{aligned}
 |\exp z|^2
   &=\exp z\,\exp\overline z
     =\exp(2\operatorname{Re}z)
     =\bigl(\exp(\operatorname{Re}z)\bigr)^2,\\
 |\exp z|&=\exp(\operatorname{Re}z).
 \end{aligned}
\]

A function $f:\mathbb C\to\mathbb C$ is continuous at $z$ if, for every
$\varepsilon>0$, there is $\delta>0$ such that
$|f(z+h)-f(z)|<\varepsilon$ whenever $|h|<\delta$.
For $|h|\leq1/2$, the geometric bound and the addition formula give
\[
 |\exp h-1|\leq\sum_{n\geq1}|h|^n\leq2|h|,\qquad
 |\exp(z+h)-\exp z|\leq2|\exp z|\,|h|.
\]
Thus the exponential is continuous.
The same definition, restricted to real inputs and outputs, defines
continuity of a real function.

For real $t$, define the cosine and sine by
\[
 \begin{aligned}
 \cos t&=\operatorname{Re}\exp(it)
       =\sum_{n\geq0}\frac{(-1)^nt^{2n}}{(2n)!},\\
 \sin t&=\operatorname{Im}\exp(it)
       =\sum_{n\geq0}\frac{(-1)^nt^{2n+1}}{(2n+1)!}.
 \end{aligned}
\]
Both functions are continuous, and $(\cos t)^2+(\sin t)^2=1$.
Conjugation gives $\exp(-it)=\cos t-i\sin t$.

A first zero of the cosine determines the required period.
At $t=2$, successive absolute terms of the cosine series decrease
from index $n=1$ onward.
Pairing the negative and positive terms after $1-2+2/3$ gives
\[
 \cos2\leq1-2+\frac23=-\frac13.
\]
For $0<t\leq2$, successive absolute terms of the sine series decrease,
since their ratio is $t^2/((2n+2)(2n+3))\leq2/3$.
Pairing the positive and negative terms after $t-t^3/6$ gives
\[
 \sin t\geq t-\frac{t^3}{6}\geq\frac t3>0.
\]
These inequalities hold first for the paired finite sums and then for
their limits.

Set
\[
 A=\{t\in[0,2]:\cos t\leq0\},\qquad \alpha_0=\inf A.
\]
Since $\cos0=1$ and $\cos2<0$, continuity gives $0<\alpha_0<2$.
If $\cos\alpha_0>0$, continuity contradicts the existence of points of $A$
arbitrarily close to $\alpha_0$ from above.
If $\cos\alpha_0<0$, continuity puts a point smaller than $\alpha_0$ in $A$.
Hence $\cos\alpha_0=0$, and $\cos t>0$ for $0\leq t<\alpha_0$.
The identity $(\cos\alpha_0)^2+(\sin\alpha_0)^2=1$ and the positive sine bound
give $\sin\alpha_0=1$.

Define $\pi=2\alpha_0$.
Then $\pi>0$, $\exp(i\pi/2)=i$, and the addition formula gives
\[
 \exp(2\pi i)=1,\qquad
 \exp(z+2\pi i)=\exp z,\qquad
 \exp(4\pi i)=1,\qquad \exp(4\pi)>1.
\]

\subsection{Period coordinates and the finite product}

The period coordinate used here is the real-linear map
\[
 \zeta_Z:L_{\mathbb R}\longrightarrow\mathbb C,
 \qquad \zeta_Z(v)=(v,e_1)+i(v,e_2)=v_1+iv_2.
\]
The lattice products are formal series in $\mathbb C[[p]]$.
For an index vector $r\in L$, its associated complex weight is
$c_r=\exp(2\pi i\zeta_Z(r))$.
Its modulus is
\[
 |c_r|=\exp\bigl(-2\pi\operatorname{Im}\zeta_Z(r)\bigr).
\]
Thus a nonzero imaginary part gives a weight of modulus different from
one. Its dependence on the chosen frame is part of its definition.

Call a lattice vector of norm two a norm-two root in this example.
The two vectors
\[
 r_1=(2,0,1),\qquad r_2=(2,0,-1)
\]
satisfy $(r_1,r_1)=(r_2,r_2)=4-2=2$.
Take the finite set $W=\{r_1,-r_1,r_2,-r_2\}$ and select
$W_+=\{r_1,r_2\}$ as one representative of each opposite pair.

An integral linear functional on $L$ is an additive map $L\to\mathbb Z$.
The functional
\[
 \nu:L\longrightarrow\mathbb Z,\qquad\nu(v)=v_1-v_3
\]
has positive values $\nu(r_1)=1$ and $\nu(r_2)=3$ on these representatives.
Set $m_{r_1}=-2$ and $m_{r_2}=1$.
Both period coordinates equal $2$, so $c_{r_1}=c_{r_2}=1$.

The product indexed by $W_+$ is therefore
\[
 \prod_{r\in W_+}(1-c_rp^{\nu(r)})^{m_r}
   =\frac{1-p^3}{(1-p)^2}.
\]
The positive exponent counts one odd generator of degree three, and the
negative exponent counts two even generators of degree one, under
Proposition~\ref{prop:pc-realization}.

Frame dependence can be computed without varying the lattice.
Use the frame
\[
 e'_1=e_2,\qquad e'_2=-e_1.
\]

Both vectors have norm one and their pairing is zero.
Their change-of-basis matrix relative to $(e_1,e_2)$ is
$\left(\begin{smallmatrix}0&-1\\1&0\end{smallmatrix}\right)$, whose
determinant is one. Thus the new frame is orthonormal and has the fixed
orientation. This change of frame is the quarter-turn used here.
Write $Z'=(P,e'_1,e'_2)$.

Substitution in the definition gives
\[
 \zeta_{Z'}(v)=(v,e_2)-i(v,e_1)=-i\zeta_Z(v).
\]
The two selected roots have period coordinate $-2i$ at $Z'$,
so their weights both become $\exp(4\pi)$.
The degree-one coefficient of the resulting product is
$2\exp(4\pi)>2$, whereas it was $2$ at $Z$.
Since both products have constant coefficient
one, even multiplication by a constant scalar cannot identify them.
Thus this construction uses the frame as actual input.

A lattice isometry is an invertible additive map $g:L\to L$ such that
$(g(v),g(w))=(v,w)$ for all $v,w\in L$.
The norm, degree, and multiplicity assignments give the displayed finite
product without supplying a transformation law under lattice isometries.

\section{Information lost by a scalar series}

Theorem~\ref{thm:pc-integrality} recovers total exponents at each degree
when all weights equal one. Labelled factorizations can still differ.
Splitting $m_r$ into two integers with the same degree and weight leaves
the product unchanged. If weights may vary, even the degrees of the
factors can change, as the polynomial identity
\[
 (1-p)(1+p)=1-p^2
\]
shows. The left side has two degree-one factors with weights $1$ and
$-1$. The right side has one degree-two factor with weight $1$.

Signed counting also forgets pairs of opposite parity.
Let
\[
 D=k[u]\otimes_k\Lambda(\eta),\qquad
 \deg u=\deg\eta=1,
\]
with $u$ even and $\eta$ odd. Then
\[
 S_D(p)=(1-p)^{-1}(1-p)=1.
\]

For each $n\geq1$, $D_n$ has even basis $u^n$ and odd basis $u^{n-1}\eta$.
Thus its signed dimension is zero although its ordinary dimension is two.
The algebra $k$ concentrated in degree zero has the same signed
dimension series. These graded spaces have different dimensions.

Even an ordinary dimension series does not determine multiplication.
Consider the graded algebra $E=k[t]/(t^3)$ with $\deg t=1$.
Compare it with the vector space
$E'=k\,1\oplus k a\oplus k b$,
where $\deg a=1$, $\deg b=2$, and every product of $a$ and $b$ is zero.
Give $1$ the identity action. These rules define an associative algebra:
a triple of positive-degree elements has product zero under either
parenthesization, and the identity handles all other cases.
Both algebras have dimension series $1+p+p^2$.

A graded algebra isomorphism would send $t$ to $\lambda a$ for some
$\lambda\ne0$. It would then send the nonzero element $t^2$ to
$\lambda^2a^2=0$, contradicting injectivity. Hence the series does not
determine the graded algebra.

There is also no differential determined by a signed dimension series.
On the already defined space $D$, a differential means a linear map
preserving degree, reversing parity, and having square zero.
The zero map $d_0$ is one such differential.
Another is the map $d_1$ defined on the basis by
\[
 d_1(u^n)=0\quad(n\geq0),\qquad
 d_1(u^{n-1}\eta)=u^n\quad(n\geq1).
\]
Every image is even and is killed by $d_1$, so $d_1^2=0$.

For a differential $d$, its cohomology in degree $n$ and parity $\epsilon$
is
\[
 H_{n,\epsilon}(D,d)=
 \frac{\ker(d:D_{n,\epsilon}\to D_{n,\epsilon+\bar1})}
      {\operatorname{im}(d:D_{n,\epsilon+\bar1}\to D_{n,\epsilon})}.
\]
The condition $d^2=0$ makes the denominator a subspace of the numerator.

For $d_0$, this cohomology is $D$ itself.
For $d_1$, the map from odd to even is an isomorphism in each degree
$n\geq1$. Both cohomology spaces therefore vanish there, and only the
even constant in degree zero remains.
The underlying monomial count is identical for these two differentials.
